% !TeX root = ../main.tex

\sysusetup{
  keywords  = {多模态RAG，文档理解，视觉定位，混合检索，溯源，强化学习，路由决策},
  keywords** = {Mémoire de bachelor, L'Abstract, Mots-clés},
}

\begin{abstract}
  本开题报告围绕复杂排版文档的检索与问答需求，提出一套多模态文档智能助手方案。系统前端进行版面分析与 OCR，抽取文本、表格、图像与视频片段等结构化元素；文本由 BGE-m3 向量化，表格经 OCR 转为 Markdown 后向量化，图像通过 ViT 等视觉编码提取特征，视频采用运动感知 MLLM 获得时序表征，构建多模态索引。在线阶段引入意图分类路由器，基于 GRPO 强化学习与评测模型打分更新策略，选择合适的 VLM（Qwen2.5-VL、LLaMA 3.2 Vision、DeepSeek-VL 等），并执行混合检索与 BGE-Reranker-v2-m3 重排序。生成阶段将检索到的文本、表格与视觉描述拼接为上下文，通过提示约束输出页码与图表来源等溯源信息。评测数据集采用 DocVQA 与 ChartQA 子集等，指标包括 Recall@K、MRR、ANLS 与 EM，并记录响应延迟。预期在单卡环境下实现可部署、可解释的多模态文档问答系统原型。

\end{abstract}

% 法语（第三种语言）摘要，本科不需要就全部注释掉即可，研究生请直接注释
%\begin{abstract**}
%  La longeur de l'abstract français doit faire référence à l'abstract chinois. Le titre de l'abstract français est «ABSTRACT». La première lettre de chaque mot-clé doit être en gros, et les mots-clés doivent être séparer par des virgules et espaces. Les mots-clés de chaque langue doivent être correspondants les uns aux autres.
%\end{abstract**}
